\documentclass[12pt,twoside,a5paper]{article}
% http://www-h.eng.cam.ac.uk/help/tpl/textprocessing/latex_maths+pix/node6.html symboles de math
% http://fr.wikibooks.org/wiki/Programmation_LaTeX Programmation latex (wikibook)
%=========================== En-Tete =================================
%--- Insertion de paquetages (optionnel) ---
\usepackage[french]{babel}   % pour dire que le texte est en fran{\'e}ais
\usepackage{a4}	             % pour la taille   
\usepackage[T1]{fontenc}     % pour les font postscript
\usepackage{epsfig}          % pour gerer les images
\usepackage{amsmath, amsthm} % tres bon mode mathematique
\usepackage{amsfonts,amssymb}% permet la definition des ensembles

\usepackage{lastpage}

\usepackage[top=1.0cm, bottom=1.0cm, left=1.0cm, right=1.0cm]{geometry}
% gauche, haut, droite, bas, entete, ente2txt, pied, txt2pied
% \usepackage{vmargin}
% \setmarginsrb{1.0cm}{1.0cm}{1.0cm}{1.0cm}{15pt}{3pt}{57pt}{3pt}

\usepackage{lscape} % changement orientation page

\pagestyle{empty}

% % % en-tete et pieds de page configurables : fancyhdr.sty
% http://www.trustonme.net/didactels/250.html
% http://ww3.ac-poitiers.fr/math/tex/pratique/entete/entete.htm
% http://www.ctan.org/tex-archive/macros/latex/contrib/fancyhdr/fancyhdr.pdf
% \usepackage{fancyhdr}
% \pagestyle{fancy}
% % \newcommand{\chaptermark}[1]{\markboth{#1}{}}
% % \newcommand{\sectionmark}[1]{\markright{\thesection\ #1}}
% \fancyhf{}
% \fancyhead[LE,RO]{\bfseries\thepage}
% \fancyhead[LO]{\bfseries\rightmark}
% \fancyhead[RE]{\bfseries\leftmark}
% \fancyfoot[LE]{\thepage /\pageref{LastPage} \hfill
	% TITLE
% \hfill \includegraphics[width=0.5cm]{img/logo_glider.png} }
% \fancyfoot[RO]{\includegraphics[width=0.5cm]{img/logo_glider.png} \hfill
	% TITLE
% \hfill \thepage /\pageref{LastPage}}
% \renewcommand{\headrulewidth}{0.5pt}
% \renewcommand{\footrulewidth}{0.5pt}
% \addtolength{\headheight}{0.5pt}
% \fancypagestyle{plain}{
	% \fancyhead{}
	% \renewcommand{\headrulewidth}{0pt}
% }

%--- Definitions de nouvelles couleurs ---
% \definecolor{verylightgrey}{rgb}{0.8,0.8,0.8}
% \definecolor{verylightgray}{gray}{0.80}
% \definecolor{lightgrey}{rgb}{0.6,0.6,0.6}
% \definecolor{lightgray}{gray}{0.6}

\usepackage{fancybox}
\usepackage{tikz}
\usepackage{graphics}

%% \newcommand{TitleVerticalOnLeft}[1][notitle]{ \begin{minipage}[h]{0.05\textwidth} \rotatebox{90}{ \textsc{#1} } \end{minipage} \hfill \begin{minipage}[h]{0.95\textwidth} }%

%============================= Corps =================================
\begin{document}

\setlength\parindent{0pt} % \noindent for all document

SUR LES FRONTIÈRES~\\~\\

%% L'Héritier/L'Héritière
\begin{minipage}[h]{0.05\textwidth} \rotatebox{90}{ \textsc{L'Héritier/L'Héritière} } \end{minipage} \hfill \begin{minipage}[h]{0.95\textwidth}
%% \TitleVerticalOnLeft{L'Héritier/L'Héritière}
	Je suis \hrulefill~le/la \hrulefill ~\\
	Qui \hrulefill ~\\
	<< \hrulefill >>~\\
	Faille :  \hrulefill~(  \hrulefill )~\\
	Ami : \hrulefill ~\\
	Ennemi : \hrulefill Âge :  \hrulefill ~\\
	Destin :  \hrulefill~\Ovalbox{~} \Ovalbox{~} \Ovalbox{~} \Ovalbox{~} \Ovalbox{~}~\\
\end{minipage}

\begin{minipage}[h]{0.45\textwidth}
	\begin{tikzpicture}[scale=0.75]
		% \draw [very thin, gray] (-5,-5) grid ( 5, 5);
		% \draw (-2, 0) node {-2, 0};
		
		\draw [black]( 0.80, 0.50) node {ÂME};
		\draw [black](-0.90, 0.50) node {CORPS};
		\draw [black]( 0.00,-1.00) node {ESPRIT};
	
		\draw [thick, black]   ( 0.00, 0.00) -- ( 1.75,-1.00)arc(-30: 90:2.00cm) -- cycle;
		\draw [thick, black]   ( 0.00, 0.00) -- ( 0.00, 2.00)arc( 90:210:2.00cm) -- cycle;
		\draw [thick, black]   ( 0.00, 0.00) -- (-1.75,-1.00)arc(210:330:2.00cm) -- cycle;
		
		\draw [thick, black]   ( 0.00, 0.00) -- ( 3.50,-2.00)arc(-30: 90:4.00cm) -- cycle;
		\draw [thick, black]   ( 0.00, 0.00) -- ( 0.00, 4.00)arc( 90:210:4.00cm) -- cycle;
		\draw [thick, black]   ( 0.00, 0.00) -- (-3.50,-2.00)arc(210:330:4.00cm) -- cycle;
	\end{tikzpicture}
\end{minipage} \hfill \begin{minipage}[h]{0.45\textwidth}

	Caractéristiques \& Notes~\newline~\newline~\newline~\newline~\newline~\newline~\newline
	
	%% Traumatismes
	\begin{minipage}[h]{0.05\textwidth} \rotatebox{90}{ \textsc{Traumatismes} } \end{minipage} \hfill \begin{minipage}[h]{0.95\textwidth}
		{ \scriptsize
			\Ovalbox{~} Humilié \emph{Pas de caractéristiques}~\newline
			\Ovalbox{~} Déchu 	\emph{Pas de don/trait/devise}~\newline
			\Ovalbox{~} Blessé 	\emph{Circonstances d{\'e}favorables}~\newline
			\Ovalbox{~} Mutilé 	\emph{Ne lance plus les d{\'e}s}~\newline
			\Ovalbox{~} Mort 	\emph{Ne peut plus agir}~\newline
		}
	\end{minipage}
	
\end{minipage}

\hrule~\\

%% Résolution~\\
\begin{minipage}[h]{0.05\textwidth} \rotatebox{90}{ \textsc{Résolution} } \end{minipage} \hfill \begin{minipage}[h]{0.95\textwidth}
	{ 	\scriptsize
		Lancez 1D6 et ajouter la caractéristique appropriée. Comptez vos succès en fonction du total :
		\begin{itemize}
			\item[] 1-2 : 0 succès / 3-5 : 1 succès / 6-8 : 2 succès / 9-11 : 3 succès / 12 : 4 succès
		\end{itemize}
		Le trait, le don et la devise donnent des circonstances favorables : vous lancez 2D6 et gardez le meilleur~\\
		La faille et certains traumatismes donnent des circonstances défavorables : vous lancez 2D6 et gardez le moins bon.~\\
		Pour réussir une action, vous devez cumuler un certain nombre de succès :
		\begin{itemize}
			\item[] Action difficile : 1 succès / ardue : 2 succès / inhumaine : 3 succès / impossible : 4 succès
		\end{itemize}
	}
\end{minipage}

\clearpage

\hrule~\\

%% \textsc{Héritier}~\\
\begin{minipage}[h]{0.05\textwidth} \rotatebox{90}{ \textsc{Héritier} } \end{minipage} \hfill \begin{minipage}[h]{0.95\textwidth}
	NOM \hrulefill~FAILLE \hrulefill~\\
	AMI \hrulefill~ENNEMI \hrulefill~\\
	DEVISE \hrulefill~\\
	DESTIN \hrulefill~\Ovalbox{~} \Ovalbox{~} \Ovalbox{~} \Ovalbox{~} \Ovalbox{~}~\\
\end{minipage}

\hrule~\\

%% \textsc{Héritier}~\\
\begin{minipage}[h]{0.05\textwidth} \rotatebox{90}{ \textsc{Héritier} } \end{minipage} \hfill \begin{minipage}[h]{0.95\textwidth}
	NOM \hrulefill~FAILLE \hrulefill~\\
	AMI \hrulefill~ENNEMI \hrulefill~\\
	DEVISE \hrulefill~\\
	DESTIN \hrulefill~\Ovalbox{~} \Ovalbox{~} \Ovalbox{~} \Ovalbox{~} \Ovalbox{~}~\\
\end{minipage}

\hrule~\\

%% \textsc{Héritier}~\\
\begin{minipage}[h]{0.05\textwidth} \rotatebox{90}{ \textsc{Héritier} } \end{minipage} \hfill \begin{minipage}[h]{0.95\textwidth}
	NOM \hrulefill~FAILLE \hrulefill~\\
	AMI \hrulefill~ENNEMI \hrulefill~\\
	DEVISE \hrulefill~\\
	DESTIN \hrulefill~\Ovalbox{~} \Ovalbox{~} \Ovalbox{~} \Ovalbox{~} \Ovalbox{~}~\\
\end{minipage}

\hrule~\\

%% \textsc{MOMENTS}
\begin{minipage}[h]{0.05\textwidth} \rotatebox{90}{ \textsc{Moments} } \end{minipage} \hfill \begin{minipage}[h]{0.95\textwidth}
	\begin{tabular}[c]{ p{0.35\textwidth}|p{0.18\textwidth}|p{0.18\textwidth}|p{0.18\textwidth}|}
		\textsc{Héritiers}				&					&					&					\\ \hline
		\textsc{Épreuve : Corps}		&					&					&					\\ \hline
		\textsc{Épreuve : Esprit}		&					&					&					\\ \hline
		\textsc{Épreuve : Âme}			&					&					&					\\ \hline
		\textsc{Question : Barbares}	&					&					&					\\ \hline
		\textsc{Question : Bêtes}		&					&					&					\\ \hline
		\textsc{Question : Héritiers}	&					&					&					\\ \hline
		\textsc{Question : Héros}		&					&					&					\\ \hline
		\textsc{Question : Empire}		&					&					&					\\ \hline
		\textsc{Question : Famille}		&					&					&					\\ \hline
		\textsc{Question : Ami}			&					&					&					\\ \hline
		\textsc{Question : Ennemi}		&					&					&					\\ \hline
		\textsc{Question : Réalité}		&					&					&					\\ \hline
	\end{tabular}~\\
\end{minipage}

\hrule

\clearpage

%% \textsc{Contrée}~\\
\begin{minipage}[h]{0.05\textwidth} \rotatebox{90}{ \textsc{Contrée} } \end{minipage} \hfill \begin{minipage}[h]{0.95\textwidth}
	\textsc{Nom de la communauté} 			\hrulefill~\\
	\textsc{Quel est le paysage ?} 			\hrulefill~\\
	\textsc{Qui est le sénéchal ?} 			\hrulefill~\\
	\textsc{Quelle est la forteresse ?} 	\hrulefill~\\
	\textsc{Quel est le trouble ?} 			\hrulefill~\\
	\textsc{Quelle est la ressource ?} 		\hrulefill~\\
	\textsc{Quel est le secret ?} 			\hrulefill~\\
	\textsc{Qui est le héros passé ?} 		\hrulefill~\\
	\textsc{Quelles sont les bêtes ?} 		\hrulefill~\\
	\textsc{Qui sont les barbares ?} 		\hrulefill~\\
	\textsc{Quelles sont les rumeurs ?} 	\hrulefill~\\
\end{minipage}

\foreach \x in {1,2,...,25}{ \hrule~\\ }%

\clearpage

\textsc{Mémo}~\\

%% \textsc{Épreuves}
\begin{minipage}[h]{0.05\textwidth} \rotatebox{90}{ \textsc{Épreuves} } \end{minipage} \hfill \begin{minipage}[h]{0.95\textwidth}
	\begin{tabular}[c]{ p{0.20\textwidth} p{0.23\textwidth} p{0.26\textwidth} p{0.23\textwidth}|}
		Résultats	&	Succès		&	Difficulté			&	Seuil		\\ 
		3			&	1 succès	&	Action difficile	&	1 succès	\\ 
		6			&	2 succès	&	Action ardue		&	2 succès	\\ 
		9			&	3 succès	&	Action inhumaine	&	3 succès	\\ 
		12			&	4 succès	&	Action impossible	&	4 succès	\\ 
	\end{tabular}~\\
\end{minipage}

L'usage de \textbf{Don}, \textbf{Trait}, \textbf{Devise} apportent des circonstances favorables.~\\
La \textbf{Faille}, les \textbf{Traumatismes} et le \textbf{sous-nombre} apportent des circonstances défavorables.~\\
Quand les circonstances sont favorables, on lance \textbf{deux dés et on garde le meilleur}. 
Quand elles sont défavorables, on lance \textbf{deux dés et on garde le moins bon}.~\\
Si il y a des circonstances favorables et défavorables, la situation reste neutre.~\\

\textsc{Utiliser le destin}~\\
%% \begin{minipage}[h]{0.05\textwidth} \rotatebox{90}{ \textsc{Utiliser le destin} } \end{minipage} \hfill \begin{minipage}[h]{0.95\textwidth}
	\begin{minipage}[h]{0.60\textwidth}
		\small
		\begin{itemize}
			\setlength{\parskip}{0pt}
			\setlength{\itemsep}{0pt plus 1pt}
			\item[$\bullet$] Réusssir une action
			\item[$\bullet$] Réussir une action dans les temps
			\item[$\bullet$] Aider un autre à réussir une action
			\item[$\bullet$] Récupérer de toutes ses blessures
			\item[$\bullet$] Emporter un affrontement
			\item[$\bullet$] Connaître une information utile
			\item[$\bullet$] Connaître une personne qui peut aider
			\item[$\bullet$] Disposer d'un objet approprié
		\end{itemize}
	\end{minipage} \hfill \begin{minipage}[h]{0.38\textwidth}
		\small
		Les héritiers disposent de \textbf{6 points de Destin}. Quand tout est dépensé, ils sont prisonniers de leur Destin, qui finira par s'accomplir. Ils perdent l'usage de leur don et sont désormais appelés \textbf{Compagnons}.
		
		Les points de Destin doivent toujours être dépensés avant de lancer les dés.
	\end{minipage}
%% \end{minipage}

~\\

%% \textsc{Traumatismes}~\\
\begin{minipage}[h]{0.05\textwidth} \rotatebox{90}{ \textsc{Traumatismes} } \end{minipage} \hfill \begin{minipage}[h]{0.95\textwidth}
	\begin{itemize}
		\item[\Ovalbox{~}] \textsc{Humilié} Vous ne pouvez plus utiliser vos traits, dons, et devises
		\item[\Ovalbox{~}] \textsc{Déchu} Vous n'ajoutez plus vos caractéristiques à vos jets de dés
		\item[\Ovalbox{~}] \textsc{Blessé} Vous avez une faille permanente, vous pouvez être mutilé
		\item[\Ovalbox{~}] \textsc{Mutilé} Vous ne lancez plus les dés, vous pouvez mourir
		\item[\Ovalbox{~}] \textsc{Mort} Vous ne pouvez plus agir. Vous pouvez toujours influencer le récit avec des points de Destin et conseiller les Héritiers survivants 
	\end{itemize}
\end{minipage}

%% \clearpage
%% Affrontements
%% Protagonistes

\end{document}
